\section{Introduction}
\label{sec:intro}

The proliferation of Parameter-Efficient Fine-Tuning (PEFT), particularly through Low-Rank Adaptation (LoRA) \cite{hu21lora}, has revolutionized the deployment of large-scale deep learning models. Rather than storing and serving multiple independent copies of a large pre-trained model for different tasks, researchers can now maintain a single frozen base model backbone (e.g., Vision Transformers or deep residual networks) and inject lightweight task-specific low-rank experts. 
When deployed in multi-task streaming environments, a serving system must dynamically route incoming requests to their respective specialized experts. This paradigm, known as dynamic model merging, has emerged as a key solution to optimize both memory footprint and serving latency.

However, executing dynamic routing in deep architectures presents a fundamental challenge known as the \emph{routing paradox}. To avoid the prohibitive computational cost of executing multiple parallel forward passes through the heavy backbone, the routing decision must be determined early in the network---often at an early representation layer ($l_{\text{route}}$). But early representation features are high-dimensional, highly entangled, and corrupted by heteroscedastic task-agnostic noise. This noise propagates downstream, leading to routing jitter and degraded ensembling accuracy.

To mitigate this noise, existing systems utilize dynamic activation-blending (such as SABLE or PAC-ZCA), where intermediate activations are dynamically interpolated layer-by-layer based on early-routing scores. While these methods are highly resilient, their calibration remains a critical bottleneck. Specifically, under ultra-low data regimes (e.g., $N=16$ calibration samples per task), optimizing layer-wise routing parameters via standard Empirical Risk Minimization (ERM) suffers from severe \emph{transductive overfitting}. Without regularization, the optimizer exploits high-frequency noise directions across deep layers, causing extreme parameter oscillations (high-frequency temperature spikes) and failing to generalize to out-of-sample test streams.

In multi-task serving environments, systems also face \emph{Vectorization Collapse}. When processing a batched stream of mixed tasks, if weight-space merging is used to dynamically average expert weights on a per-request basis, the GPU cannot batch these requests together because their model weights differ. The serving framework is forced to run sequential, single-sequence forward passes, which destroys GPU tensor parallelism and collapses vectorized throughput. Activation-blending (PAC-STM) completely avoids Vectorization Collapse because the base backbone weights remain 100\% frozen and shared. This allows the framework to process heterogeneous requests in a single vectorized batch, applying only cheap sample-specific scaling of low-rank adapter activation outputs.

\begin{figure}[t]
\centering
\includegraphics[width=\linewidth]{results/fig2.png}
\caption{\textbf{Layer-wise Temperature Trajectories.} Comparison of learned log-temperatures across network layers. Unregularized Temp-Only ERM oscillates wildly due to transductive overfitting on tiny calibration sets ($N=16$). Our proposed PAC-STM utilizes the derived Markovian trajectory prior to enforce smooth depth-wise transition, matching the Oracle behavior.}
\label{fig:trajectories}
\end{figure}

This paper addresses this bottleneck from a rigorous learning-theoretic perspective. We reject the dominant heuristic paradigm of model merging, arguing that empirical success on a specific stream is fragile without solid mathematical guarantees explaining \emph{why} the learned ensembling parameters generalize. We present **PAC-Bayesian Smooth Trajectory Merging (PAC-STM)**, a framework that models the routing parameters across deep network layers as a probability distribution over trajectories. 

By formulating the prior over these trajectories as an autoregressive Markov chain (a Gaussian random walk across depth), we prove that the Kullback-Leibler (KL) complexity penalty in the PAC-Bayesian bound analytically derives a first-order finite-difference smoothness regularizer. This bridges the gap between randomized generalization bounds and stable, deterministic trajectory optimization, restricting the capacity of the deep routing head without sacrificing layer-specific expressiveness.

We make the following core contributions:
\begin{itemize}
    \item \textbf{Unit-Norm PCA Subspace Projection (UN-PCA-SEP):} We propose a preprocessing step that projects normalized hidden representations onto task-specific principal component bases, bounding coordinate energies strictly in $[0, 1]$ and eliminating coordinate magnitude ambiguity under extreme noise.
    \item \textbf{Markovian Trajectory Prior:} We formulate a Gaussian random walk prior over layer-wise log-temperatures, capturing depth-wise correlations and reflecting the physical continuity of representation learning in deep neural networks.
    \item \textbf{Analytical Trajectory KL-Regularizer:} We prove a closed-form theorem for the KL-divergence of our Markovian posterior, providing an exact, parameter-free smoothness penalty for trajectory optimization.
    \item \textbf{Rigorous Empirical Validation:} We evaluate PAC-STM on a simulated 14-layer Analytical Coordinate Sandbox (ICS) designed to model representation manifolds under strict noise controls. PAC-STM mitigates transductive overfitting, outperforming unregularized ERM by up to $2.05\%$ on heterogeneous batching while maintaining complete, simulation-validated immunity to both \emph{Heterogeneity Collapse} and \emph{Vectorization Collapse} on mixed serving streams.
\end{itemize}
